\documentclass[12pt,a4paper]{article}

% ─── Packages ───────────────────────────────────────────────────────────────
\usepackage{amsmath, amssymb, amsthm}
\usepackage{geometry}
\usepackage{booktabs}
\usepackage{array}
\usepackage{xcolor}
\usepackage{titlesec}
\usepackage{hyperref}
\usepackage{graphicx}
\usepackage{mathtools}
\usepackage{bm}           % bold math
\usepackage{microtype}
\usepackage{parskip}

\geometry{margin=2.5cm}

% ─── Colours ────────────────────────────────────────────────────────────────
\definecolor{accent}{RGB}{30, 80, 160}
\definecolor{eqbg}{RGB}{245, 248, 255}

% ─── Section formatting ─────────────────────────────────────────────────────
\titleformat{\section}{\large\bfseries\color{accent}}{}{0em}{}[\titlerule]
\titleformat{\subsection}{\normalsize\bfseries\color{accent!80!black}}{}{0em}{}

% ─── Custom environments ────────────────────────────────────────────────────
\newtheorem*{defn}{Definition}
\newcommand{\eq}[1]{\colorbox{eqbg}{\ensuremath{\displaystyle #1}}}

% ─── Macros ─────────────────────────────────────────────────────────────────
\newcommand{\vv}[1]{\bm{#1}}          % bold vector
\newcommand{\R}{\mathbb{R}}
\newcommand{\T}{\hat{\bm{T}}}         % unit tangent
\newcommand{\N}{\hat{\bm{N}}}         % unit normal
\newcommand{\B}{\hat{\bm{B}}}         % unit binormal

% ════════════════════════════════════════════════════════════════════════════
\begin{document}

\begin{center}
  {\LARGE \textbf{Mathematical Equations for the Parametric Bird}}\\[6pt]
  {\large \textit{A Complete Formulation of the Wireframe Bird Image}}\\[4pt]
  {\small Generated by: \texttt{bird.py} \quad|\quad Surface Model v7}
\end{center}

\vspace{0.5em}
\hrule
\vspace{1em}

% ────────────────────────────────────────────────────────────────────────────
\section{Master Surface Equation}

The entire bird is assembled from \textbf{tube-swept parametric surfaces}.
Every component curve $\vv{\gamma}(t)$ generates a surface by sweeping a
circle of radius $r(t)$ perpendicular to the curve:

\begin{equation}
  \boxed{
    \vv{S}(t,\theta)
    \;=\;
    \vv{\gamma}(t)
    \;+\;
    r(t)\!\left[\,\hat{\vv{N}}(t)\cos\theta
                  +\, \hat{\vv{B}}(t)\sin\theta\,\right]
  }
  \qquad t\in[0,1],\quad \theta\in[0,2\pi]
\end{equation}

where $\hat{\vv{N}}(t)$ and $\hat{\vv{B}}(t)$ are the \emph{normal} and
\emph{binormal} vectors of the \textbf{Rotation-Minimising Frame} (RMF)
described in Section~\ref{sec:rmf}.

% ────────────────────────────────────────────────────────────────────────────
\section{Bézier Curve Primitives}

All spine curves $\vv{\gamma}(t)$ are constructed from Bézier curves over
control points $\vv{P}_i \in \R^3$.

\subsection{Quadratic Bézier Curve}
\begin{equation}
  \vv{B}_2(t;\,\vv{P}_0,\vv{P}_1,\vv{P}_2)
  \;=\;
  (1-t)^2\,\vv{P}_0
  + 2(1-t)t\,\vv{P}_1
  + t^2\,\vv{P}_2,
  \qquad t\in[0,1]
\end{equation}

\subsection{Cubic Bézier Curve}
\begin{equation}
  \vv{B}_3(t;\,\vv{P}_0,\vv{P}_1,\vv{P}_2,\vv{P}_3)
  \;=\;
  (1-t)^3\,\vv{P}_0
  + 3(1-t)^2t\,\vv{P}_1
  + 3(1-t)t^2\,\vv{P}_2
  + t^3\,\vv{P}_3,
  \qquad t\in[0,1]
\end{equation}

% ────────────────────────────────────────────────────────────────────────────
\section{Rotation-Minimising Frame (RMF)}
\label{sec:rmf}

The standard Frenet–Serret frame is undefined where curvature vanishes.
We use the \textbf{Rotation-Minimising Frame} (Bishop frame), propagated
iteratively:

\paragraph{Step 1 — Unit Tangent.}
\begin{equation}
  \T(t_i)
  \;=\;
  \frac{\dfrac{d\vv{\gamma}}{dt}\bigg|_{t_i}}
       {\left\lVert \dfrac{d\vv{\gamma}}{dt}\bigg|_{t_i} \right\rVert}
\end{equation}
Computed numerically via central differences
$\;\dfrac{d\vv{\gamma}}{dt}\big|_{t_i} \approx
 \dfrac{\vv{\gamma}(t_{i+1})-\vv{\gamma}(t_{i-1})}{t_{i+1}-t_{i-1}}$.

\paragraph{Step 2 — Initial Normal.}
Choose an arbitrary reference $\vv{e}$
(e.g.\ $\vv{e}=(0,1,0)$) not parallel to $\T(t_0)$:
\begin{equation}
  \hat{\vv{N}}(t_0)
  \;=\;
  \frac{\T(t_0)\times\vv{e}}{\lVert\T(t_0)\times\vv{e}\rVert},
  \qquad
  \hat{\vv{B}}(t_0) = \T(t_0)\times\hat{\vv{N}}(t_0)
\end{equation}

\paragraph{Step 3 — Parallel Transport (RMF propagation).}
For $i = 1, 2, \ldots, n-1$:
\begin{align}
  \vv{N}^\perp_i
  &\;=\;
  \hat{\vv{N}}(t_{i-1})
  - \bigl[\hat{\vv{N}}(t_{i-1})\cdot\T(t_i)\bigr]\,\T(t_i)
  \\[4pt]
  \hat{\vv{N}}(t_i)
  &\;=\;
  \frac{\vv{N}^\perp_i}{\lVert\vv{N}^\perp_i\rVert}
  \\[4pt]
  \hat{\vv{B}}(t_i)
  &\;=\;
  \frac{\T(t_i)\times\hat{\vv{N}}(t_i)}
       {\lVert\T(t_i)\times\hat{\vv{N}}(t_i)\rVert}
\end{align}

% ────────────────────────────────────────────────────────────────────────────
\section{Wing Surface Equations}

Each wing is modelled as a \textbf{bilinear ruled surface} stretched between
a \emph{leading-edge} curve $\vv{\alpha}(u)$ and a \emph{trailing-edge}
curve $\vv{\beta}(u)$, both cubic Bézier curves in $\R^3$.

\subsection{Ruled Wing Surface}
\begin{equation}
  \boxed{
    \vv{S}_{\mathrm{wing}}(u,v)
    \;=\;
    (1-v)\,\vv{\alpha}(u) \;+\; v\,\vv{\beta}(u),
    \qquad u,v\in[0,1]
  }
\end{equation}

\begin{center}
\begin{tabular}{ll}
  $u=0$ & wing root (body junction) \\
  $u=1$ & wing tip \\
  $v=0$ & leading edge \\
  $v=1$ & trailing edge
\end{tabular}
\end{center}

\subsection{Left Wing Leading Edge}
\begin{equation}
  \vv{\alpha}_L(u)
  = \vv{B}_3\!\left(u;\;
    \begin{pmatrix}0.00\\0.10\\0.10\end{pmatrix},\;
    \begin{pmatrix}-1.80\\0.05\\0.55\end{pmatrix},\;
    \begin{pmatrix}-3.80\\-0.08\\0.60\end{pmatrix},\;
    \begin{pmatrix}-5.50\\-0.20\\0.15\end{pmatrix}
  \right)
\end{equation}

\subsection{Left Wing Trailing Edge}
\begin{equation}
  \vv{\beta}_L(u)
  = \vv{B}_3\!\left(u;\;
    \begin{pmatrix}0.00\\0.10\\-0.22\end{pmatrix},\;
    \begin{pmatrix}-1.40\\0.00\\-0.05\end{pmatrix},\;
    \begin{pmatrix}-3.60\\-0.12\\0.25\end{pmatrix},\;
    \begin{pmatrix}-5.40\\-0.20\\0.12\end{pmatrix}
  \right)
\end{equation}

\subsection{Right Wing (Bilateral Symmetry)}
The right wing is the reflection of the left wing across the sagittal plane
($x = 0$):
\begin{equation}
  \vv{\alpha}_R(u) = \vv{M}\,\vv{\alpha}_L(u),
  \qquad
  \vv{\beta}_R(u)  = \vv{M}\,\vv{\beta}_L(u),
  \qquad
  \vv{M} = \mathrm{diag}(-1,\;1,\;1)
\end{equation}

\subsection{Wing Mesh Lines}
The wing surface is rendered as two families of iso-parameter lines:

\begin{align}
  \text{Span lines (chord-dir):}\quad
  &\mathcal{L}_u \;=\; \bigl\{\vv{S}(u_0,v)\;:\;v\in[0,1]\bigr\}
  \quad\text{for fixed }u_0 \\[4pt]
  \text{Chord lines (span-dir):}\quad
  &\mathcal{L}_v \;=\; \bigl\{\vv{S}(u,v_0)\;:\;u\in[0,1]\bigr\}
  \quad\text{for fixed }v_0
\end{align}

% ────────────────────────────────────────────────────────────────────────────
\section{Body / Torso}

The torso is a \textbf{swept tube} along a cubic Bézier spine from head to
tail:

\subsection{Body Spine Curve}
\begin{equation}
  \vv{\gamma}_{\mathrm{body}}(t)
  = \vv{B}_3\!\left(t;\;
    \begin{pmatrix}0\\0\\0.90\end{pmatrix},\;
    \begin{pmatrix}0\\0.05\\0.50\end{pmatrix},\;
    \begin{pmatrix}0\\0\\-0.15\end{pmatrix},\;
    \begin{pmatrix}0\\-0.05\\-0.70\end{pmatrix}
  \right)
\end{equation}

\subsection{Body Radius Profile}

The radius follows an \textbf{asymmetric sine-power} profile, thicker at
the breast and tapering at head and tail:
\begin{equation}
  r_{\mathrm{body}}(t)
  \;=\;
  R_0
  + A\,\sin^{\alpha}\!\bigl(\pi\,t\bigr),
  \qquad t\in[0,1]
\end{equation}

\begin{center}
\begin{tabular}{lll}
  \toprule
  Parameter & Symbol & Value \\
  \midrule
  Base radius      & $R_0$    & $0.06$ \\
  Amplitude        & $A$      & $0.28$ \\
  Shape exponent   & $\alpha$ & $1.2$  \\
  \bottomrule
\end{tabular}
\end{center}

The body surface is then:
\begin{equation}
  \vv{S}_{\mathrm{body}}(t,\theta)
  =
  \vv{\gamma}_{\mathrm{body}}(t)
  + r_{\mathrm{body}}(t)\,
  \bigl[\hat{\vv{N}}(t)\cos\theta + \hat{\vv{B}}(t)\sin\theta\bigr]
\end{equation}

% ────────────────────────────────────────────────────────────────────────────
\section{Head Spike}

The beak/head is a \textbf{tapered cone} swept along a straight line above
the body:

\subsection{Spike Axis}
\begin{equation}
  \vv{\gamma}_{\mathrm{spike}}(t)
  = (1-t)\,\vv{p}_{\mathrm{head}} + t\,\vv{p}_{\mathrm{tip}},
  \qquad t\in[0,1]
\end{equation}

where $\vv{p}_{\mathrm{head}} = \vv{\gamma}_{\mathrm{body}}(0)$ and
$\vv{p}_{\mathrm{tip}} = \vv{p}_{\mathrm{head}} + (0,\,0,\,0.45)^\top$.

\subsection{Spike Radius Profile}

The radius tapers to zero at the tip via a power law:
\begin{equation}
  r_{\mathrm{spike}}(t)
  \;=\;
  R_{\mathrm{spike}}\,(1-t)^{\,p},
  \qquad
  R_{\mathrm{spike}} = 0.10,\quad p = 1.4
\end{equation}

% ────────────────────────────────────────────────────────────────────────────
\section{Tail Feathers}

$M$ tail feathers fan out below the body in a symmetric spread.
Each feather $j=0,\ldots,M{-}1$ is a swept tube along a
\textbf{quadratic Bézier curve}.

\subsection{Azimuth Angle}
\begin{equation}
  \phi_j
  = \phi_{\min}
  + \frac{j}{M-1}\,(\phi_{\max} - \phi_{\min}),
  \qquad
  \phi_{\min} = 200^\circ,\quad
  \phi_{\max} = 340^\circ
\end{equation}

\subsection{Feather Length (Bell Curve)}
\begin{equation}
  L_j
  = L_{\mathrm{side}}
  + \bigl(L_{\mathrm{mid}} - L_{\mathrm{side}}\bigr)
    \sin\!\left(\pi\,\frac{j}{M-1}\right),
  \qquad
  L_{\mathrm{mid}} = 2.2,\quad L_{\mathrm{side}} = 0.80
\end{equation}

\subsection{Feather Spine (Quadratic Bézier)}
\begin{align}
  \vv{p}_{\mathrm{tip},j}
  &= \vv{p}_{\mathrm{tail}}
     + L_j\,
     \begin{pmatrix}\cos\phi_j \\ 0 \\ \sin\phi_j\end{pmatrix}
  \\[6pt]
  \vv{p}_{\mathrm{ctrl},j}
  &= \tfrac{1}{2}\!\left(\vv{p}_{\mathrm{tail}}+\vv{p}_{\mathrm{tip},j}\right)
     + \begin{pmatrix}
         0.10\,L_j\cos(\phi_j+\tfrac{\pi}{2})\\
         0.04\,L_j \\
         0.10\,L_j\sin(\phi_j+\tfrac{\pi}{2})
       \end{pmatrix}
  \\[6pt]
  \vv{\gamma}_j(t)
  &= \vv{B}_2\!\left(t;\;
       \vv{p}_{\mathrm{tail}},\;
       \vv{p}_{\mathrm{ctrl},j},\;
       \vv{p}_{\mathrm{tip},j}
     \right)
\end{align}

\subsection{Feather Radius Profile}
\begin{equation}
  r_j(t)
  \;=\;
  R_{0,T}\,(1-t)^{\,\beta_T}
  + R_{\min,T},
  \qquad
  R_{0,T} = 0.18,\quad
  \beta_T  = 0.55,\quad
  R_{\min,T} = 0.008
\end{equation}

% ────────────────────────────────────────────────────────────────────────────
\section{Summary of All Equations}

\begin{center}
\renewcommand{\arraystretch}{1.6}
\begin{tabular}{>{\bfseries}l l}
  \toprule
  Component       & Governing Equation \\
  \midrule
  Master surface  &
    $\vv{S}(t,\theta) = \vv{\gamma}(t) + r(t)\bigl[\N\cos\theta + \B\sin\theta\bigr]$ \\
  Quadratic Bézier &
    $\vv{B}_2(t) = (1-t)^2\vv{P}_0 + 2(1-t)t\,\vv{P}_1 + t^2\vv{P}_2$ \\
  Cubic Bézier    &
    $\vv{B}_3(t) = (1-t)^3\vv{P}_0 + 3(1-t)^2t\,\vv{P}_1 + 3(1-t)t^2\vv{P}_2 + t^3\vv{P}_3$ \\
  RMF transport   &
    $\N_i = \dfrac{\N_{i-1} - (\N_{i-1}\cdot\T_i)\T_i}{\lVert\cdots\rVert}$,\;
    $\B_i = \T_i\times\N_i$ \\
  Wing surface    &
    $\vv{S}_{\mathrm{wing}}(u,v) = (1-v)\vv{\alpha}(u) + v\,\vv{\beta}(u)$ \\
  Body radius     &
    $r_{\mathrm{body}}(t) = R_0 + A\sin^{\alpha}(\pi t)$ \\
  Spike radius    &
    $r_{\mathrm{spike}}(t) = R_{\mathrm{spike}}\,(1-t)^{1.4}$ \\
  Tail azimuth    &
    $\phi_j = \phi_{\min} + \tfrac{j}{M-1}(\phi_{\max}-\phi_{\min})$ \\
  Tail length     &
    $L_j = L_{\mathrm{side}} + (L_{\mathrm{mid}}-L_{\mathrm{side}})\sin\!\left(\tfrac{\pi j}{M-1}\right)$ \\
  Tail radius     &
    $r_j(t) = R_{0,T}(1-t)^{\beta_T} + R_{\min,T}$ \\
  Right-wing sym. &
    $\vv{\alpha}_R(u) = \mathrm{diag}(-1,1,1)\,\vv{\alpha}_L(u)$ \\
  \bottomrule
\end{tabular}
\end{center}

% ────────────────────────────────────────────────────────────────────────────
\section{Rendering Parameters}

\begin{center}
\renewcommand{\arraystretch}{1.4}
\begin{tabular}{llr}
  \toprule
  Parameter                 & Symbol          & Value  \\
  \midrule
  Span-wise sample lines    & $N_{\mathrm{span}}$  & 110 \\
  Chord-wise sample lines   & $N_{\mathrm{chord}}$ & 75  \\
  Points per line           & $N_{\mathrm{pts}}$   & 200 \\
  Body ring count           & $N_{\mathrm{body}}$  & 140 \\
  Rings per cross-section   & $N_\theta$           & 28--38 \\
  Tail feather count        & $M$                  & 18  \\
  Tail sample count         & $N_{\mathrm{tail}}$  & 120 \\
  Line width (wings)        & $\ell_w$             & 0.25 pt \\
  Camera elevation          & $\epsilon$           & $18^\circ$ \\
  Camera azimuth            & $\psi$               & $-90^\circ$ \\
  \bottomrule
\end{tabular}
\end{center}

\vspace{1em}
\begin{center}
  \small\textit{All coordinates are in arbitrary units; the bird spans
  approximately $[-5.5,\,5.5]$ in the $x$-axis (wing span).}
\end{center}

\end{document}
